%!TEX root = TheoreticalPhysicsC.tex
\section{Lecture: ``Quantisation''}

How does one identify the right Hilbert space, Hamiltonian and observables for a given system? Unfortunately there is no systematic method for determining this information for a physical system. That means that, given a classical description of a system, there is no mathematically well-defined map which automatically assigns to it the correct quantum system. The reason for this is the following. Classical physics represents a good approximation to the physics of quantum mechanical systems which are exposed to interactions with their surroundings. These damaging effects of the environment on a quantum system are known as decoherence. Decoherence is the reason why classical mechanics is such a good approximation to everyday physics. Closed systems, after all, hardly exist at all: every charged object couples to the electromagnetic field, which is itself a quantum system. If we move such an object back and forth it radiates, and the radiation carries the information off to infinity at the speed of light.

Ideally one would like a way of uncovering the original quantum system from its classical mechanical description alone. That is, we want a way of undoing decoherence. Unfortunately this is impossible in general. Such a fictitious inversion procedure has become known as ``quantisation''.

\begin{greybox}[An inverse problem]
	Decoherence maps different quantum systems onto one and the same classical system -- two points land on top of each other. Nor can a face be reconstructed from the noisy image of a surveillance camera.
\end{greybox}

We evade the difficulties of decoherence by trial and error. Decades of experience have handed us a list of Hilbert spaces, Hamiltonians and observables which we now know to give the correct physical description of broad classes of systems. Indeed, all of the fundamental interactions (with the exception of gravity) have been understood in this way. The aim of this lecture is to draw up a list of the basic quantum systems known to model certain idealised systems correctly. Never think that there is a simple recipe for it -- if it were that easy we would need no experimental physics.


\subsection{A particle in one dimension}
The most important system to have been successfully modelled with quantum mechanics is the particle in one dimension. As always we dutifully work our way through all five postulates. The Hilbert space (Postulate 1) is in this case $\mathcal{H}\equiv L^2(\mathbb{R})$. That is, we describe preparations (Postulate 2) of pure states of a particle with the help of complex functions -- \emph{wavefunctions} -- $\Psi:\mathbb{R}\rightarrow \mathbb{C}$. Wavefunctions are written in ket notation as
\begin{equation}
	|\Psi\rangle \equiv \int_{-\infty}^\infty \Psi(x) |x\rangle \, dx
\end{equation}
The restriction to pure states is no arbitrary choice: every density matrix is a mixture of pure states. Only normalised vectors are allowed, $\langle \Psi|\Psi\rangle = 1$.

So far the wavefunction is nothing but notation; it acquires meaning only through measurement. A theory needs both: states \emph{and} detectors.

According to Postulate 3, detectors are modelled by operators. Here we describe position detectors only. Such detectors deliver ``yes''/``no'' results, i.e.\ either the particle is at a given location or it is not. In order to describe such detectors we need an interval $[a,b]\subset \mathbb{R}$, where $a<b$. The corresponding detector is defined by the operator
\begin{equation}
	F_{[a,b]} \equiv \int_{a}^b |x\rangle\langle x|\, dx
\end{equation}
Such operators do not cover all physically possible detectors, but they do serve as a physically relevant foundation. Why a region and not a point? Because the question of the exact location always has the answer ``probability zero''. We are not tied to intervals: more complicated regions are taken care of by additivity -- for disjoint sets $S_1$, $S_2$ we have $F(S_1\cup S_2) = F(S_1)+F(S_2)$.

Suppose that the system is in a pure state $|\Psi\rangle$ and that we are about to carry out the measurement $F_{[a,b]}$. With what probability will we find the particle in the region $a \le x \le b$? Postulate 4 gives us the response rate of the detector $F_{[a,b]}$ for a preparation $\rho \equiv |\Psi\rangle\langle \Psi|$ as
\begin{equation}
	\begin{split}
		\mathbb{P}[\text{``$x\in [a,b]$''}] &= \tr(\rho F_{[a,b]}) \\
		&= \tr (|\Psi\rangle\langle \Psi|F_{[a,b]})\\
		&= \tr (\langle \Psi|F_{[a,b]}|\Psi\rangle) \\
		&= \langle \Psi|F_{[a,b]}|\Psi\rangle \\
		&= \int_{-\infty}^\infty dx \int_{-\infty}^\infty dy\, \Psi^*(x)\Psi(y) \langle x|F_{[a,b]}|y\rangle   \\
		&= \int_{-\infty}^\infty dx \int_{-\infty}^\infty dy \int_{a}^b dz\, \Psi^*(x)\Psi(y) \langle x|z\rangle\langle z|y\rangle \\
		&= \int_{-\infty}^\infty dx \int_{-\infty}^\infty dy \int_{a}^b dz\, \Psi^*(x)\Psi(y) \delta(x-z)\delta(z-y) \\
		&= \int_{a}^b dz\, \Psi^*(z)\Psi(z) \\
		&= \int_{a}^b |\Psi(z)|^2 \, dz.
	\end{split}
\end{equation}
In the case that $a\rightarrow -\infty$ and $b\rightarrow \infty$ we obtain the condition
\begin{equation}
	1 = \mathbb{P}[\text{``$x\in \mathbb{R}$''}].
\end{equation} 
Since the quantity $|\Psi(x)|^2$ is non-negative, we may interpret the quantity
\begin{equation}
	\eqbox{p(x) \equiv |\Psi(x)|^2}
\end{equation}
as a probability density.

\begin{intuition}
	The wavefunction assigns a \emph{complex} number to every possible position and delivers the position only indirectly: the probability of finding the particle between $a$ and $b$ is the area under the curve $p(x)$. You should resist the temptation to picture $\Psi$ itself as a density.
\end{intuition}

% fig_vl07_dichte.tex — Warum $\Psi$ keine Dichte ist und $|\Psi|^2$ eine:
% dieselbe Wellenfunktion einmal als Real-/Imaginärteil und einmal als
% Wahrscheinlichkeitsdichte, deren Fläche über einem Intervall die
% Ansprechwahrscheinlichkeit des Detektors $F_{[a,b]}$ ist.
%
% Folgt den Konventionen von fig_vl11_kasten.tex.
%
% Geometrie aus der Formel, nicht freihand.  Beide Panels zeigen ein und
% denselben normierten Zustand — ein Gaußsches Wellenpaket, den einfachsten
% Zustand mit $\Psi(x)\ne\delta(x-y)$, den die Vorlesung fordert:
%
%     Psi(x) = (2 pi sigma^2)^(-1/4) exp(-(x-x0)^2/(4 sigma^2)) exp(i k x),
%
% mit sigma = 1, x0 = 0, k = 3.  Daraus folgt alles Gezeichnete:
%
%   Re Psi  = A exp(-x^2/4) cos(3x),   A = (2 pi)^(-1/4) = 0.63162,
%   Im Psi  = A exp(-x^2/4) sin(3x),
%   |Psi|   = A exp(-x^2/4)                      (die Einhüllende),
%   p(x)    = |Psi(x)|^2 = exp(-x^2/2)/sqrt(2 pi),  Maximum 0.39894.
%
% Die Trägerwelle exp(ikx) hebt sich im Betragsquadrat exakt weg: die
% Oszillation in (a) hat in (b) keine Spur hinterlassen.  Genau das ist
% in Prosa mühsam und im Bild eine Sekunde.
%
% Die eingefärbte Fläche ist keine Schätzung, sondern
%     P[x in [a,b]] = Phi(1.5) - Phi(-0.5) = 0.62466,
% mit a = -0.5, b = 1.5 und der Standardnormalverteilung Phi.
% Der schmale Streifen bei c = -1.8 hat Breite eps = 0.2 und damit
% Fläche p(c) eps = 0.0158: dieselbe Konstruktion, nur schmaler.  Er ist
% der Grund, warum F_{x} = |x><x| kein Operator ist — schrumpft eps auf
% null, schrumpft die Fläche mit.
%
% Entwurfsentscheidungen:
%   * Farbe trägt genau zwei Bedeutungen.  In (a) unterscheidet sie die
%     beiden reellen Komponenten desselben komplexen Objekts (Blau/Orange);
%     in (b) ist alles, was Wahrscheinlichkeit ist, blau — der breite und
%     der schmale Streifen unterscheiden sich nur in der Sättigung, weil
%     sie dieselbe Größe sind, nur verschieden breit.
%   * Referenzen sind neutral grau: die Einhüllende +-|Psi| in (a) ist
%     gepunktet grau, keine Daten.
%   * Die beiden Wahrscheinlichkeiten sind gefüllt und nicht schraffiert:
%     schraffiert wird hier, was verboten ist; hier ist die Fläche selbst
%     die Größe, die man ablesen soll, und eine Schraffur würde ihren
%     Inhalt (die Integralformel) unlesbar machen.
%   * Im Sinne des Graustufendrucks ist Im Psi gestrichelt und Re Psi
%     durchgezogen; die beiden Flächen unterscheiden sich zusätzlich
%     durch ihre Breite, nicht nur durch den Grauwert.
%   * Beide Panels benutzen denselben x-Bereich und denselben Maßstab,
%     damit die Breiten vergleichbar sind: die Einhüllende von Psi ist eine
%     Gaußkurve der Breite sqrt(2) sigma, die Dichte p eine der Breite
%     sigma.  Das Quadrieren verschmälert die Glocke um genau diesen
%     Faktor sqrt(2) — im Bild sofort zu sehen, in Prosa eine Fußnote.
%   * Der Streifen sitzt bei c = -1.8 und nicht dichter am Intervall,
%     damit er nicht als dessen Rand missverstanden wird.

\begin{center}
\begin{tikzpicture}

% ─────────────────────────────────────────────────────────────
%  (a)  Die Wellenfunktion selbst: komplex, vorzeichenbehaftet
% ─────────────────────────────────────────────────────────────
\begin{axis}[
    name=psi,
    tpcplot,
    width=6.6cm, height=6.2cm,
    grid=none,
    axis x line=middle,
    axis y line=none,
    axis line style={-{Stealth[length=5pt]}, spacecadet, line width=0.6pt},
    xtick=\empty,
    xmin=-3.6, xmax=3.9,
    ymin=-0.94, ymax=0.94,
    clip=false,
    domain=-3.6:3.75, samples=401,
    title={(a) $\Psi$ \tpcifenglish{is complex}{ist komplex}},
]

  % Einhüllende +-|Psi| = +-A exp(-x^2/4): Referenz, daher grau
  \addplot[birrengray!85, line width=0.7pt, densely dotted]
    { 0.63162*exp(-x^2/4)};
  \addplot[birrengray!85, line width=0.7pt, densely dotted]
    {-0.63162*exp(-x^2/4)};

  % Real- und Imaginärteil derselben Wellenfunktion
  \addplot[wavefn, cgblue]
    {0.63162*exp(-x^2/4)*cos(deg(3*x))};
  \addplot[birrenorange, line width=1.1pt, densely dashed]
    {0.63162*exp(-x^2/4)*sin(deg(3*x))};

  % Direkte Beschriftung statt Legende: jeweils am eigenen Extremum.
  % Re Psi hat sein Maximum A bei x = 0, Im Psi sein Minimum -0.58978
  % bei x = -pi/6 = -0.5236.
  \node[lbl, text=cgblue, anchor=south]      at (axis cs:0,0.66)       {$\real\Psi$};
  \node[lbl, text=birrenorange, anchor=north] at (axis cs:-0.524,-0.68) {$\imag\Psi$};

  % Einhüllende ansprechen, ohne sie zu berühren
  \node[lbl, text=birrengray, font=\scriptsize, anchor=west]
    at (axis cs:-3.55,0.44) {$\pm|\Psi|$};
  \draw[birrengray!85, line width=0.4pt]
    (axis cs:-2.90,0.42) -- (axis cs:-2.45,0.155);

  \node[lbl, text=birrengray, font=\scriptsize, anchor=east]
    at (axis cs:-3.62,0.02) {$0$};

  % Achsenbeschriftung als Knoten unter der Pfeilspitze, nicht als xlabel:
  % ein xlabel ragt rechts aus dem Rahmen und stößt im 12pt-Satz gegen
  % die y-Beschriftung von (b).
  \node[lbl, anchor=north] at (axis cs:3.55,-0.06) {$x$};
\end{axis}

% ─────────────────────────────────────────────────────────────
%  (b)  Das Betragsquadrat: Fläche = Wahrscheinlichkeit
% ─────────────────────────────────────────────────────────────
\begin{axis}[
    tpcplot,
    at={(psi.south east)}, anchor=south west, xshift=22mm,
    width=6.6cm, height=6.2cm,
    grid=none,
    ylabel={$p(x)$},
    every axis y label/.style={at={(axis description cs:-0.17,0.5)},
                               rotate=90, anchor=south, font=\small},
    xmin=-3.6, xmax=3.9,
    ymin=0, ymax=0.52,
    xtick={-1.8,-0.5,1.5}, xticklabels={$c$,$a$,$b$},
    ytick={0,0.1,0.2,0.3,0.4},
    clip=false,
    domain=-3.6:3.75, samples=301,
    title={(b) $p=|\Psi|^{2}$ \tpcifenglish{is the density}{ist die Dichte}},
]

  % P[x in [a,b]] = Flaeche unter p zwischen a und b
  \addplot[draw=none, fill=cgblue!20, domain=-0.5:1.5, samples=101]
    {exp(-x^2/2)/sqrt(2*pi)} \closedcycle;
  \draw[cgblue!55, line width=0.5pt] (axis cs:-0.5,0) -- (axis cs:-0.5,0.35207);
  \draw[cgblue!55, line width=0.5pt] (axis cs: 1.5,0) -- (axis cs: 1.5,0.12952);

  % Derselbe Schnitt, nur um eps = 0.2 breit: gleiche Farbe, dunkler,
  % weil es dieselbe Größe ist.
  \addplot[draw=none, fill=cgblue!60, domain=-1.9:-1.7, samples=11]
    {exp(-x^2/2)/sqrt(2*pi)} \closedcycle;

  % Die Dichte
  \addplot[spacecadet, line width=1.1pt] {exp(-x^2/2)/sqrt(2*pi)};

  % Was die Fläche ist, steht in der Fläche
  \node[lbl, anchor=center] at (axis cs:0.5,0.15) {$\int_a^b p\,dx$};

  % Die Fläche beziffert: die Aussage steht über der Fläche, die sie meint
  \node[lbl, anchor=center] at (axis cs:0.5,0.455)
    {$\mathbb{P}\bigl[x\in[a,b]\bigr]=\tpcifenglish{0.62}{0{,}62}$};

  \node[lbl, anchor=north] at (axis cs:3.55,-0.012) {$x$};

  % Der schmale Streifen und seine Pointe
  \node[lbl, text=spacecadet, font=\scriptsize, anchor=south]
    at (axis cs:-1.8,0.095) {$\varepsilon$};
  \node[lbl, font=\scriptsize, anchor=west, align=left] at (axis cs:-3.40,0.275)
    {$\varepsilon\to0$\\[-1pt]$\Rightarrow\mathbb{P}\to0$};
  \draw[birrengray!85, line width=0.4pt]
    (axis cs:-2.14,0.225) -- (axis cs:-1.90,0.130);
\end{axis}

\end{tikzpicture}
\end{center}


On the left of the figure $\Psi$ swings through zero and carries a sign; on the right only the smooth bell $p$ remains. Watch the narrow strip: as its width shrinks, the area shrinks with it.

Warning: the object
\begin{equation}
	F_{\{x\}} \equiv |x\rangle \langle x|,
\end{equation}
does \emph{not} represent an operator. This has to do with the interpretation of $p(x) \, (\equiv \tr(F_{\{x\}}\rho))$: the probability of finding the particle \emph{exactly} at the point $x$ is always equal to zero. One must always specify an interval in order to obtain a non-vanishing probability.

Mathematically one sees the same thing in the failure of $F_{\{x\}}$ to have the projector property $F^2 = F$:
\begin{equation}
	F_{\{x\}}^2 = |x\rangle \langle x|x\rangle \langle x| = \langle x|x\rangle\, |x\rangle\langle x| = \delta(0)\, F_{\{x\}},
\end{equation}
and $\delta(0)$ is by definition infinite: instead of $F_{\{x\}}$ we get $\infty\cdot F_{\{x\}}$. For an interval, by contrast, what sits in the middle is $\langle x|y\rangle = \delta(x-y)$:
\begin{equation}
	\begin{split}
		F_{[a,b]}^2 &= \int_a^b \int_a^b |x\rangle \langle x|y\rangle \langle y|\, dx\, dy \\
		&= \int_a^b \int_a^b |x\rangle\, \delta(x-y)\, \langle y|\, dx\, dy \\
		&= \int_a^b |x\rangle\langle x|\, dx = F_{[a,b]}.
	\end{split}
\end{equation}
The delta function disposes of one of the integrals, and what is left over is the operator we started with: $F_{[a,b]}$ really is a projector.
\begin{merksatz}
	Stand-alone ket-bra combinations of the position basis are to be avoided -- they are not operators. Under an integral the very same expressions are harmless.
\end{merksatz}
All the same, one sometimes works with such ``projectors'' on a purely informal level. In particular it is helpful to picture the ``projector'' $|x\rangle\langle x|$ as representing the event that the particle sits at the position $x$. In this manner one introduces the position operator (following Eq.~(\ref{eq:observableA}))
\begin{equation}
	\eqbox{\widehat{x} \equiv \int_{-\infty}^{\infty} x |x\rangle \langle x|\, dx}
\end{equation}
We read it as the \emph{first moment} of the position measurement; the hat is there to remind us that this is an operator and not a number. It acts as a multiplication operator,
\begin{equation}
	\widehat{x}|\Psi\rangle = \int_{-\infty}^{\infty} x\, |x\rangle \langle x|\Psi\rangle\, dx = \int_{-\infty}^{\infty} x\, \Psi(x)\, |x\rangle\, dx,
\end{equation}
so $\Psi(x)$ turns into $x\,\Psi(x)$.

An important remark: it is physically impossible to prepare the particle in the ``classical state'' $|\Psi\rangle = |x\rangle$. You would need a device of infinite precision. The search for eigenvectors tells us the same thing: a $|u\rangle$ with $\widehat{x}|u\rangle = \lambda|u\rangle$ would have a delta function as its wavefunction and would therefore be an element of the position basis. One nevertheless speaks of ``position eigenvectors'' all the time -- they are simply idealised objects, not elements of $\mathcal{H}$. We will come back to this topic in a later lecture. Naturally occurring pure states are in general always in superposition:
\begin{equation}
	\Psi(x) \not= \delta(x-y).
\end{equation}


The expectation value of $\widehat{x}$ represents the average position of the particle:
\begin{equation}
	\begin{split}
		\langle \widehat{x}\rangle &= \tr(\widehat{x}\rho) \\
		&= \int x|\Psi(x)|^2\, dx,\\
		&= \int x p(x)\, dx.
	\end{split}
\end{equation}
(This equation holds only in the case that the particle is in the pure state $\rho \equiv |\Psi\rangle\langle \Psi|$.) Behind the symbol $\langle \widehat{x}\rangle$ there lies a whole procedure: prepare, measure, write down -- many times over, until the average settles.

The momentum operator of a one-dimensional quantum mechanical particle is defined by
\begin{equation}
	\eqbox{\widehat{p} \equiv -i\hbar \frac{d}{dx}}
\end{equation}
Why? One can indeed put forward many heuristic arguments for this identification, but the honest answer is that nature simply wanted it this way. How does $\widehat{p}$ act on a state $|\Psi\rangle$? This action is given by the following ``matrix elements'' (these hold for all $|\Psi\rangle \in \mathcal{H}$):
\begin{equation}
	\langle x|\widehat{p}|\Psi\rangle \equiv -i\hbar \frac{d}{dx}\Psi(x).
\end{equation}

Postulate 5: the time evolution of a particle is given in terms of the Hamiltonian (corresponding to the \emph{kinetic energy})
\begin{equation}
	\widehat{H} \equiv \frac{\widehat{p}^2}{2m} = -\frac{\hbar^2}{2m}\frac{d^2}{dx^2}
\end{equation}
with matrix elements
\begin{equation}
	\langle x|\widehat{H}|\Psi\rangle = -\frac{\hbar^2}{2m}\frac{d^2}{dx^2}\Psi(x)
\end{equation}
Here $m$ denotes the mass of the particle. Where does this expression come from? From the only quantisation rule we have: the Hamiltonian is the classical energy, with a hat on it. Classically $E = p^2/2m$; we put the hat on. This is a rule of thumb, not a procedure -- every quantisation remains a guess, and only experiment can decide.

\begin{example}
	Suppose a particle has been prepared in an \emph{equal superposition} over an interval $[-1,1]$: the corresponding state/wavefunction reads
	\begin{equation}
		\Psi(x) \equiv \begin{cases}
			\frac{1}{\sqrt2}, \quad x\in [-1,1] \\
			0, \quad \text{otherwise}.
		\end{cases}
	\end{equation}
	Before we calculate, let us guess: by symmetry the mean value should be exactly zero. The normalisation $\tr(\rho) = \tr(|\Psi\rangle\langle \Psi|) = \langle \Psi|\Psi\rangle$ of the state is given by
	\begin{equation}
		\begin{split}
			\langle \Psi|\Psi\rangle &= \int_{-\infty}^\infty |\Psi(x)|^2 \, dx \\
			&= \int_{-1}^1  \frac{1}{2} \, dx \\
			&= 1
		\end{split}
	\end{equation} 
	The expectation value of the position operator $\widehat{x}$ is calculated as follows:
	\begin{equation}
		\begin{split}
			\langle \widehat{x}\rangle = \tr(\widehat{x}\rho) &= \langle \Psi|\widehat{x}|\Psi\rangle \\
		&= \int_{-1}^1 x \times \frac{1}{2}\, dx = \frac12\left[\frac{x^2}{2}\right]_{x= -1}^1 = 0.
		\end{split}
	\end{equation}
\end{example}
